% -------------------------------------------------------
%  Section 6: Reported results
% -------------------------------------------------------
\section{نتایج گزارش‌شده در مقاله}
\label{sec:eval}

جدول ۵ مقاله ۲۰ پیکربندی را گزارش می‌کند (شش برای \lr{VGG16} و چهارده برای \lr{ResNet50}) و جدول \ref{tab:papermain} چکیده آن را بر حسب خانواده طرح می‌دهد. نخستین مشاهده این است که ستون صحت آموزش برای همه ۲۰ ردیف میان \pct{۹۹٫۳۰} و \pct{۱۰۰} است؛ یعنی هیچ قدرت تفکیکی ندارد و تمام بار مقایسه روی یک عدد واحد، یعنی زیان لگاریتمی، می‌افتد.

\begin{table}[!ht]
    \centering
    \footnotesize
    \caption{چکیده نتایج گزارش‌شده مقاله بر حسب خانواده طرح (زیان لگاریتمی، نمره خصوصی؛ کمتر بهتر است)}
    \label{tab:papermain}
    \begin{tabular}{|l|c|c|c|c|}
    \hline
    \textbf{خانواده طرح} & \textbf{تعداد پیکربندی} & \textbf{بهترین} & \textbf{بدترین} & \textbf{میانگین} \\ \hline
    \lr{S1}--\lr{S3} (تصویر خام)          & ۶ & \num{۰٫۰۲۶۵} & \num{۰٫۰۵۴۷} & \num{۰٫۰۳۷۸} \\ \hline
    \lr{S4} (\lr{Split})                  & ۶ & \num{۰٫۰۳۲۰} & \num{۰٫۰۵۸۷} & \num{۰٫۰۴۴۷} \\ \hline
    \lr{S5} همراه تصویر خام               & ۵ & \num{۰٫۰۲۷۹} & \num{۰٫۰۵۲۱} & \num{۰٫۰۳۵۳} \\ \hline
    \lr{S5} بدون تصویر خام                & ۳ & \num{۰٫۰۴۴۶} & \num{۰٫۰۸۸۹} & \num{۰٫۰۶۱۳} \\ \hline
    \end{tabular}
\end{table}

مقاله چهار نتیجه از این جدول بیرون می‌کشد. نخست، \textbf{عرض ثابت بهتر است}: روی \lr{ResNet50} ترتیب \lr{S3} با \num{۰٫۰۲۶۵}، سپس \lr{S1} با \num{۰٫۰۳۱۶} و سپس \lr{S2} با \num{۰٫۰۳۳۰} برقرار است و \lr{S2} که بیشترین تنوع عرض را دارد روی هر دو معماری بدترین طرح خام است. دوم، \textbf{قطعه‌بندی \lr{Split} کافی است}: روی \lr{ResNet50}، \lr{S4} کارایی «مشابه» \lr{S1}--\lr{S3} دارد، پس حتی با نادیده گرفتن کامل چیدمان فضایی هم می‌توان به نتایج خوب رسید. سوم، \textbf{\lr{Mask} تنها با تصویر خام کار می‌کند}: \lr{S5} به‌تنهایی به‌دلیل تُنُکی ضعیف است ولی ترکیب آن با تصویر خام عرض‌ثابت به \num{۰٫۰۲۷۹} می‌رسد. چهارم، \textbf{بخش‌ها سهم نابرابر دارند}: \kd{.text} مهم‌ترین بخش است اما حتی بدون آن هم کارایی قابل قبول می‌ماند، و \kd{.rsrc} ویژگی‌های دسته‌بندی بسیار خوبی نشان می‌دهد که مقاله آن را به اشتراک منابعی مانند آیکون و منو میان گونه‌های یک خانواده نسبت می‌دهد.

نکته مهم آن است که بهترین ردیف جدول از آنِ \lr{S3} است و هیچ پیکربندی قطعه‌بندی از آن جلو نمی‌زند؛ یعنی ادعای مقاله «بهترین صحت» نیست، بلکه رسیدن به نتیجه‌ای \emph{نزدیک} در ازای استقلال از چیدمان فضایی است. این تمایز در \S\ref{sub:transfer} تعیین‌کننده خواهد بود.
